\documentclass[11pt]{article}
\usepackage[T1]{fontenc}
\usepackage{lmodern,microtype}
\usepackage[margin=1in]{geometry}
\usepackage{amsmath,amssymb,amsthm,mathtools}
\usepackage{booktabs,tabularx,array,enumitem}
\usepackage[dvipsnames]{xcolor}
\usepackage[colorlinks=true,linkcolor=MidnightBlue,citecolor=MidnightBlue,urlcolor=MidnightBlue]{hyperref}
\usepackage{bookmark}
\setlength{\emergencystretch}{2em}
\setcounter{tocdepth}{1}
\setlist{topsep=4pt,itemsep=3pt,parsep=0pt}
\newtheorem{theorem}{Theorem}[section]
\newtheorem{proposition}[theorem]{Proposition}
\newtheorem{lemma}[theorem]{Lemma}
\newtheorem{corollary}[theorem]{Corollary}
\theoremstyle{definition}
\newtheorem{definition}[theorem]{Definition}
\theoremstyle{remark}
\newtheorem{remark}[theorem]{Remark}
\numberwithin{equation}{section}
\newcommand{\R}{\mathbb R}
\newcommand{\C}{\mathbb C}
\newcommand{\dd}{\,\mathrm d}
\newcommand{\ee}{\mathrm e}
\newcommand{\ii}{\mathrm i}
\newcommand{\calJ}{\mathcal J}
\newcommand{\abs}[1]{\left|#1\right|}
\newcommand{\norm}[1]{\left\lVert#1\right\rVert}
\DeclareMathOperator{\atanh}{atanh}
\newcolumntype{Y}{>{\raggedright\arraybackslash}X}
\renewcommand{\arraystretch}{1.12}
\title{Green Potentials and Boundary Logarithms\\for the Newman Family}
\author{Drafted for Edward Baker}
\date{8 November 2026\\\small Working analytical draft}
\hypersetup{pdftitle={Green Potentials and Boundary Logarithms for the Newman Family},pdfauthor={Drafted for Edward Baker},pdfsubject={Complete Green identities, scoped approximation obstructions, and logarithmic acquisition}}
\begin{document}
\maketitle
\begin{abstract}
We develop a logarithmic approach to detecting nonreal zeros of the
canonical Newman family. Starting from the complete theta source, we
derive the logarithmic resolvent and an exact positive strip Green
potential, including boundary zeros, multiplicities and quantitative
tails. At time zero the potential becomes a critical-line logarithmic
average minus an absolutely convergent prime baseline, with the zeta
pole correction retained. Two results delimit possible estimating
methods. A fixed translated-source family has common source and count
bounds, but an order-\(\eta\) source and safe-line resolvent change
produces an order-\(\sqrt\eta\) Green change. A separate theorem,
using the short-mollifier moment bound of Soundararajan as presented
by Radziwi{\l}{\l}, gives a positive uniform average floor for a
specified class of quadratic certificates. We derive the normalized
transfer and incorporate the corrected real mixed-term calculation.
We then prove a one-dimensional sublevel estimate and an explicit
one-sided logarithmic approximation bound through real multiple zeros.
It permits zeros of the approximant and requires no lower modulus on
the real boundary. The local exponent deteriorates with height.
The complete finite-source logarithmic assembly and a favorable
canonical value remain outstanding. All results are analytical;
no numerical experiment is used as proof, and no zero exclusion is
claimed.
\end{abstract}
\tableofcontents
\clearpage
\section{The estimating problem}\label{sec:introduction}

The zeros of an entire function can be represented either by poles of
its logarithmic derivative or by point masses in the Laplacian of its
logarithmic modulus. For the Newman family these descriptions lead to
a positive, exponentially localized Green potential. They also isolate
a distinction that governs the analysis: an exact representation of a
zero signal, an accurate approximation to that signal, and an upper
bound excluding a zero are three different statements.

The field studied here is
\[
 H_\lambda(z)=\int_0^\infty \ee^{\lambda u^2}\Phi(u)\cos(zu)\dd u,
 \qquad
 H_0(z)=\frac18\xi\left(\frac12-\frac{\ii z}{2}\right),
\]
with the literal complete theta source \(\Phi\) specified in
Section~\ref{sec:source}. We use \(0\le\lambda\le1/2\) for the Green
geometry and the inherited canonical zero strip \(|\Im z_j|\le1\).
The arithmetic estimates are restricted to \(\lambda=0\).
Real symmetry means that a nonreal zero has an upper-half-plane
partner. All zero sums count multiplicity.

The main observable \(J_\lambda(X)\), defined in
Section~\ref{sec:green}, has the exact form
\[
 J_\lambda(X)=\sum_{\Im z_j>0}
 \atanh\frac{\sin(\pi\Im z_j/8)}
               {\cosh(\pi(X-\Re z_j)/8)}.
\]
Thus a zero with \(\Im z_j\ge d>0\) and
\(|\Re z_j-X|\le r\) forces an explicit positive margin. The desired
estimate has the opposite direction: a separately obtained upper
value must lie strictly below that margin. The potential does not
detect the multiplicity of a zero already on the real boundary.

Our development has four stages. First we establish the source,
resolvent, complete Green representation and exact prime subtraction.
Second we prove a quantitative stability obstruction on a specified
enlarged source class. Third we analyze quadratic approximation to a
pole-cancelled zeta field and prove, for fixed \(0<\theta<1/2\),
the residual floor
\[
 \liminf_{T\to\infty}\inf_A
 \frac1{|I_T|}\int_{I_T}\epsilon_A(X)^2\dd X
 \ge\frac{1/(2\theta)-11/12}{\zeta(5/2)^4}>0.
\]
Here a single bounded-coefficient Dirichlet polynomial of length at
most \((T/2)^\theta\) is used throughout a padded interval \(I_T\).
Finally we pass directly to the boundary logarithm, deriving a
one-sided error modulus from a safe interior anchor and a disk
growth bound. This last theorem allows real multiple zeros and does
not replace the measured logarithm by a harmonic or quadratic
surrogate.

The complete source and Green background are developed in the parent
manuscript and its assessed Green note \cite{OGParent,GreenStripProof}.
The source-class obstruction and quadratic floor are the branch
results \cite{OGConditioning,OGMollifier}, read with the correction and
assessment \cite{OGCorrection,OGAssessment}. The local logarithmic
theorem is \cite{OGBoundaryLog}. We give their derivations here in a
single notation. Classical complex-analysis ingredients and the
external mean-value theorem are attributed where used. We make no
literature-priority claim.

The two obstruction theorems have different domains. The first varies
the source away from theta and concerns stability of prescribed
observations. The second keeps the exact canonical field and varies
a restricted preconditioner. Neither gives a lower bound for the
canonical Green potential. The new local logarithmic theorem is an
acquisition result: it controls the error of retaining a signed
quantity whose favorable value is still unknown. Section~\ref{sec:tasks}
states precisely what remains to complete that passage.


\section{The complete theta source and logarithmic resolvent}\label{sec:source}

\subsection{Reflection, the endpoint and the canonical normalization}

Let
\[
 \vartheta(t)=\sum_{n\in\mathbb Z}\ee^{-\pi n^2t},\qquad
 p(u)=\frac{\ee^u}{2}\{\vartheta(\ee^{4u})-1\}
     =\ee^u\sum_{n\ge1}\ee^{-\pi n^2\ee^{4u}}.
\]
Jacobi's transformation
\(\vartheta(t)=t^{-1/2}\vartheta(1/t)\)
\cite{DLMFThetaModular} gives
\begin{equation}\label{eq:primitive-reflection}
 p(-u)-p(u)=\sinh u,\qquad p'(0)=-\frac12.
\end{equation}
Define
\begin{equation}\label{eq:source}
 \Phi(u)=\frac{p''(u)-p(u)}8
 =\ee^u\sum_{n\ge1}(2r_n(u)^2-3r_n(u))\ee^{-r_n(u)},
 \qquad r_n(u)=\pi n^2\ee^{4u}.
\end{equation}
Applying \(\partial_u^2-1\) to \eqref{eq:primitive-reflection}
proves that \(\Phi\) is even. The complete series is smooth, and
each summand is positive when \(u\ge0\), since \(r_n\ge\pi>3/2\).
Thus \(\Phi>0\) on \(\R\). Reflection is a property of the complete
sum; an evenly reflected finite prefix need not have matching odd
derivatives at zero.

With \(W_\lambda(u)=\ee^{\lambda u^2}\Phi(u)\), put
\begin{equation}\label{eq:H}
 H_\lambda(z)=\frac12\int_\R W_\lambda(u)\ee^{\ii zu}\dd u
            =\int_0^\infty W_\lambda(u)\cos(zu)\dd u.
\end{equation}
There is no division by \(H_\lambda(0)\).
For \(J^{\rm prim}(z)=\int_0^\infty p(u)\cos(zu)\dd u\),
two integrations by parts, using the endpoint in
\eqref{eq:primitive-reflection}, give
\begin{equation}\label{eq:primitive-H}
 8H_0(z)=\frac12-(1+z^2)J^{\rm prim}(z).
\end{equation}
For completeness, splitting the theta Mellin integral at \(t=1\)
and using Jacobi reflection gives
\[
 \xi(s)=\frac12+\frac{s(s-1)}2
 \int_1^\infty \frac{\vartheta(t)-1}{2}
       \{t^{s/2-1}+t^{(1-s)/2-1}\}\dd t.
\]
The integral is entire in \(s\). Substituting \(t=\ee^{4u}\)
and then \(s=(1-\ii z)/2\) turns the integral into
\(8J^{\rm prim}(z)\), while \(4s(s-1)=-(1+z^2)\).
Equation~\eqref{eq:primitive-H} therefore proves
\begin{equation}\label{eq:xi}
 H_0(z)=\frac18\xi(s(z)),\quad
 s(z)=\frac12-\frac{\ii z}2,\quad
 \xi(s)=\frac12s(s-1)\pi^{-s/2}\Gamma(s/2)\zeta(s).
\end{equation}
The reflected argument gives the same field
\cite{DLMFGreenPrime}. The constant \(1/2\) in
\eqref{eq:primitive-H} is indispensable.

\subsection{Complete domination, heat evolution and the product}

Set \(\rho(u)=\pi\ee^{4|u|}\), and recursively define
\[
 P_0(r)=2r^2-3r,\qquad
 P_{j+1}(r)=(1-4r)P_j(r)+4rP_j'(r).
\]
Write \(a(P)\) for the sum of the absolute coefficients of a polynomial,
and put
\[
 C_d(N)=\sum_{n>N}n^{2d}\ee^{-\pi(n^2-1)},\qquad
 B_j=a(P_j)C_{j+2}(0).
\]
Each \(C_d(N)\) is finite and tends to zero with \(N\).

\begin{lemma}[Complete source bound]\label{lem:source}
For each fixed integer \(j\ge0\),
\[
 |\Phi^{(j)}(u)|\le B_j\ee^{|u|}\rho(u)^{j+2}\ee^{-\rho(u)}.
\]
On \(u\ge0\), the differentiated arithmetic tail satisfies the same
bound with \(B_j\) replaced by \(a(P_j)C_{j+2}(N)\).
\end{lemma}
\begin{proof}
Differentiating a summand gives
\(\ee^uP_j(n^2\rho)\ee^{-n^2\rho}\).
Since \(\deg P_j=j+2\) and \(\rho\ge\pi>1\),
\[
 |P_j(n^2\rho)|\le a(P_j)n^{2j+4}\rho^{j+2},\qquad
 \ee^{-n^2\rho}\le\ee^{-\rho}\ee^{-\pi(n^2-1)}.
\]
Summation proves the half-line estimate and differentiated local
convergence. Complete smooth reflection gives the full-line estimate.
\end{proof}

For fixed \(L,b\ge0\) and integers \(k,d\ge0\), the integral
\[
 \int_\R |u|^k\ee^{Lu^2+(b+1)|u|}
             \rho(u)^d\ee^{-\rho(u)}\dd u
\]
is finite: substitution \(t=\pi\ee^{4|u|}\) leaves a product of
powers of \(t,\log(t/\pi)\), and
\(\exp\{(L/16)\log^2(t/\pi)-t\}\).
Consequently \(H\) is jointly entire in \((\lambda,z)\), and all
fixed derivatives on compact parameter sets pass under the integral.
In particular,
\begin{equation}\label{eq:heat}
 \partial_\lambda H_\lambda=-\partial_z^2H_\lambda.
\end{equation}
The opposite clock \(h(\tau,z)=H_{\lambda_0-\tau}(z)\) satisfies
\(h_\tau=h_{zz}\). No uniform assertion in a growing derivative
order follows from these fixed-order bounds.

We retain the canonical strip input
\begin{equation}\label{eq:zero-strip}
 |\Im z_j(\lambda)|\le1\qquad(0\le\lambda\le1/2)
\end{equation}
from the parent development \cite{OGParent,GreenStripProof}.
This is strip control, not an all-real-zero assumption.
The source bound also gives, uniformly in this compact time range,
\(\log\max_{|z|\le r}|H_\lambda(z)|=O(r\log(r+2))\).
One can obtain this bound directly by splitting the source integral
at a constant multiple of \(\log(r+2)\); beyond that point the
double exponential dominates \(r|u|+u^2/2+O(|u|)\).
Positivity gives a common lower bound for \(H_\lambda(0)\).
Jensen's formula on a disk of radius \(2r\) therefore gives a
full origin-centered count
\begin{equation}\label{eq:count}
 N_\lambda(r)\le Cr\log(r+2),\qquad r\ge1.
\end{equation}
The even Hadamard product is consequently
\begin{equation}\label{eq:product}
 H_\lambda(z)=H_\lambda(0)
 \prod_{\zeta\in\mathcal R_\lambda}(1-z^2/\zeta^2).
\end{equation}
Here \(\mathcal R_\lambda\) contains one representative of each
opposite pair, with multiplicity. The paired product converges
locally uniformly because \(\sum|\zeta|^{-2}<\infty\). The order
bound permits at most a linear exponential factor; evenness removes
its linear term, and evaluation at zero fixes the constant.
Thus no unspecified exponential factor remains.

\subsection{The logarithmic field and its denominator}

Let \(\nu_\lambda=\sum_j\delta_{z_j(\lambda)}\), with multiplicities,
and use \(\partial_z=(\partial_x-\ii\partial_y)/2\).
Local factorization \(H=(z-z_j)^m h\), with \(h\ne0\), proves
\begin{equation}\label{eq:distributions}
 \Delta\log|H_\lambda|=2\pi\nu_\lambda,\qquad
 R_\lambda=\frac{H_\lambda'}{H_\lambda}
           =2\partial_z\log|H_\lambda|,\qquad
 \partial_{\bar z}R_\lambda=\pi\nu_\lambda .
\end{equation}
These are distributional identities; \(\log|z|\) and \(1/z\) are
locally integrable, with the stated delta masses obtained by integrating
around a small circle. A multiplicity-\(m\) zero gives residue \(m\).
Differentiating \eqref{eq:product} away from its zeros gives
\[
 R_\lambda(z)=\sum_{\zeta\in\mathcal R_\lambda}
                    \frac{2z}{z^2-\zeta^2}.
\]
Pairing is part of the convergence prescription.

On zero-free space-time patches, logarithmic differentiation of
\eqref{eq:heat} yields
\begin{equation}\label{eq:burgers}
 \partial_\lambda R=-R_{zz}-2RR_z.
\end{equation}
This local equation is not assigned an additional unproved global
meaning at its poles. Likewise, solving only the first or last
distributional equation in \eqref{eq:distributions} leaves a harmonic
or holomorphic summand. The complete field and its product fix it.

Source approximation alone does not control division. If
\(|H-\widehat H|\le\varepsilon\),
\(|H'-\widehat H'|\le\varepsilon_1\), and
\(b=|\widehat H|-\varepsilon>0\), then
\begin{equation}\label{eq:ratio-error}
 \left|R-\frac{\widehat H'}{\widehat H}\right|
 \le\frac{\varepsilon_1}{b}
       +\frac{|\widehat H'|\varepsilon}{b|\widehat H|}.
\end{equation}
This follows by adding and subtracting
\(\widehat H'/H\). The positive denominator margin is a separate
requirement. It fails on a real interval containing a zero, which
motivates the integrated logarithmic method developed later.


\section{The complete strip Green potential}\label{sec:green}

\subsection{The finite-factor identity}

Fix
\begin{equation}\label{eq:kernels}
 a=\frac\pi8,\qquad K(t)=\frac1{16\cosh(at)},\qquad
 G(t,y)=\atanh\frac{\sin(ay)}{\cosh(at)}.
\end{equation}
Each copy of \(K\) has mass \(1/2\). Let \(\mu_X\) be the probability
measure on the two lines \(y=0,8\) that has density \(K(x-X)\) on
each. For an admissible function \(f\), write
\[
 \calJ[f](X)=\int\log|f|\dd\mu_X-\log|f(X+4\ii)|,
 \qquad J_\lambda(X)=\calJ[H_\lambda](X).
\]

The Green function of \(0<\Im z<B\), normalized by
\(-\Delta G_B=2\pi\delta\), is
\begin{equation}\label{eq:strip-green}
 G_B(z,\zeta)=\frac12\log
 \frac{\cosh(\pi(\Re z-\Re\zeta)/B)-\cos(\pi(\Im z+\Im\zeta)/B)}
      {\cosh(\pi(\Re z-\Re\zeta)/B)-\cos(\pi(\Im z-\Im\zeta)/B)}.
\end{equation}
Indeed \(z\mapsto\ee^{\pi z/B}\) maps the strip to the upper
half-plane; the logarithm of the reflected-distance ratio gives
\eqref{eq:strip-green}. Its singularity is \(-\log|z-\zeta|\).
At \(B=8\), observation \(X+4\ii\), its value is
\(G(X-\Re\zeta,\Im\zeta)\). Boundary normal derivatives divided
by \(2\pi\) give exactly the two densities \(K\).

For an interior zero \(\zeta\), the singularities of
\(\log|z-\zeta|+G_B(z,\zeta)\) cancel. This harmonic function has the
linear factor's boundary logarithm. Truncating the strip vertically
proves its Poisson representation: the side harmonic measures decay
exponentially while these logarithms grow only logarithmically.
It follows that \(\calJ[z-\zeta]\) is the positive Green value.
An exterior zero gives zero. A boundary zero is handled by its
locally integrable logarithm, or by a limit from the exterior; it
gives no half-weight atom. Finite multiplicities cause no change.

Apply this identity to every complete opposite pair in a finite
truncation of \eqref{eq:product}. Scalar logarithms cancel because
the boundary masses add to one. The only remaining issue is passage
to the complete product.

\subsection{Absolute logarithmic convergence}

Choose
\[
 m_0=\inf_\lambda H_\lambda(0)>0,\qquad
 M_1\ge\sup_{\lambda,|z|\le1}|H_\lambda'(z)|,\quad M_1>0,
\]
and set
\[
 r_*=\min(1,m_0/(2M_1)),\quad
 C_0=\sup_\lambda|\log H_\lambda(0)|,\quad
 C_*=C_0+20C(1+|\log r_*|).
\]
The segment integral of \(H'\) shows that no zero lies in
\(|z|\le r_*\). All constants are uniform for \(0\le\lambda\le1/2\).

\begin{lemma}[Unit-interval logarithmic bound]\label{lem:unit-log}
For \(v\in\R\), \(0\le b\le8\), and \(Q=|v|+12\),
\begin{equation}\label{eq:unit-log}
 \int_v^{v+1}|\log|H_\lambda(x+\ii b)||\dd x
 \le C_*Q\log^2(3Q).
\end{equation}
Moreover the series of paired factor logarithms is absolutely
convergent in local \(L^1\).
\end{lemma}
\begin{proof}
On the interval, \(|z|\le Q\). For a representative \(|\zeta|\le2Q\),
split the logarithm into those of \(z-\zeta\), \(z+\zeta\), and
\(-2\log|\zeta|\). Each positive distance logarithm is at most
\(\log(3Q)\). Its negative part has integral at most two, since
\(|x+\ii b-\zeta|\ge|x-\Re\zeta|\) and
\(\int_{-1}^1-\log|t|\dd t=2\).
The factor therefore costs at most
\(4\log(3Q)+4+2|\log r_*|\).
There are at most \(CQ\log(2Q+2)\) such representatives, by
\eqref{eq:count} and opposite pairing.

For \(|\zeta|>2Q\),
\[
 \left|\log\left|1-\frac{z^2}{\zeta^2}\right|\right|
 \le\frac43\frac{Q^2}{|\zeta|^2}.
\]
Partial summation with half the full count bounds their sum by
\[
 \frac23CQ\{\log(2Q+2)+1\}.
\]
These bounds, together with \(C_0\), give \eqref{eq:unit-log}.
They bound the sum of absolute factor logarithms, which proves the
stronger local \(L^1\) assertion.
\end{proof}

\begin{theorem}[Complete positive potential]\label{thm:green}
For every real \(X\) and \(0\le\lambda\le1/2\), both boundary
logarithms are absolutely integrable against \(K(x-X)\), and
\begin{equation}\label{eq:green-identity}
 J_\lambda(X)=
 \int_\R K(x-X)\{\log|H_\lambda(x)|+\log|H_\lambda(x+8\ii)|\}\dd x
 -\log|H_\lambda(X+4\ii)|
 =\sum_{j:y_j>0}G(X-x_j,y_j).
\end{equation}
The upper-zero sum is absolutely convergent, with all multiplicities.
There is no omitted harmonic term. The potential is nonnegative and
unchanged by multiplication of the field by a nonzero scalar.
\end{theorem}
\begin{proof}
The polynomial-logarithmic bound \eqref{eq:unit-log} is summable
against the exponentially decaying Poisson weights. Thus the finite
paired-product identity passes to the boundary integrals in \(L^1\).
At \(X+4\ii\), the product converges away from every zero by
\eqref{eq:zero-strip}. This proves the complete identity. Its finite
Green sums are nonnegative and increase to the displayed finite
value, hence converge absolutely. The product has already fixed all
harmonic data. Scalar invariance follows directly from the total
boundary mass one.
\end{proof}

\subsection{Two different tails and the exclusion margin}

For \(R\ge0\), let \(n=\lfloor R\rfloor\), \(q_0=\ee^{-a}\), and
\(S=|X|+13+n\). The explicit two-boundary estimate is
\begin{align}\label{eq:boundary-tail}
 &\int_{|x-X|>R}K(x-X)
  \{|\log|H_\lambda(x)||+|\log|H_\lambda(x+8\ii)||\}\dd x
 \nonumber\\
 &\hspace{1cm}\le
 \frac{C_*}{2}\ee^{-an}S\log^2(3S)
 \frac{1+4q_0+q_0^2}{(1-q_0)^4}.
\end{align}
To derive it, cover both sides by unit shells and use
\(K(t)\le\ee^{-a|t|}/8\) and Lemma~\ref{lem:unit-log} on both
boundaries. The remaining sum is bounded by
\(S\log^2(3S)\sum_{j\ge0}q_0^j(1+j)^3\).
The generating function of this cubic moment is the rational factor
in \eqref{eq:boundary-tail}.

For the sharper zero tail we use the additional canonical unit-count
input from the parent analysis \cite{SmoothLocalizationProof}:
after a fixed enlargement of \(A\),
\[
 \#\{j:u\le x_j(\lambda)\le u+1\}\le A\log(2+|u|)
\]
uniformly for \(0\le\lambda\le1/2\).
This inherited input includes the endpoint \(\lambda=0\); it is
separate from the global Jensen count proved above. Since
\begin{equation}\label{eq:G-decay}
 0<G(t,y)\le K_Gy\ee^{-a|t|},\qquad
 K_G=\frac{2a}{\cos^2a},\quad 0<y\le1,
\end{equation}
unit shells give
\begin{equation}\label{eq:zero-tail}
 \sum_{\substack{y_j>0\\|x_j-X|>R}}G(X-x_j,y_j)
 \le2AK_G\ee^{-an}
 \left\{\frac{\log(3+|X|)+\log(n+1)}{1-q_0}
       +\frac{q_0}{(n+1)(1-q_0)^2}\right\}.
\end{equation}
Here \(\atanh u\le u/(1-u^2)\), \(\sin(ay)\le ay\), and
\(\operatorname{sech}(at)\le2\ee^{-a|t|}\) prove
\eqref{eq:G-decay}. For the shell sum use
\(\log(3+|X|+k)\le\log(3+|X|)+\log(k+1)\) and
\(\log(n+j+1)\le\log(n+1)+j/(n+1)\).
The boundary and zero tails pay different truncations. In particular
\(J_\lambda(X)\ll_A\log(3+|X|)\), but that coarse bound is insufficient
for the following consumer.

\begin{corollary}[A single nonreal-zero margin]\label{cor:margin}
If \(d\le y_j\le1\) and \(|x_j-X|\le r\), then
\begin{equation}\label{eq:margin}
 J_\lambda(X)\ge m(d,r):=\frac{\sin(a)d}{\cosh(ar)}.
\end{equation}
An independently acquired strict upper bound below \(m(d,r)\)
excludes this specified zero class.
\end{corollary}
\begin{proof}
Use positivity of every term, \(\atanh u\ge u\), and the concavity
bound \(\sin(ay)\ge y\sin a\) for \(0\le y\le1\).
\end{proof}

An all-real slice has \(J=0\); any upper zero makes \(J>0\) at every
center. Accordingly universal time-zero vanishing is an RH-equivalent
reformulation, not an additional estimate. Uniform tails on compact
center intervals and continuity of finite root multisets show that
\(J\) is continuous through collisions. A pair approaching a double
real zero contributes a quantity linear in its upper height; when
that height has square-root time scale, so does the Green signal.
Continuity alone does not imply Lipschitz continuity or monotonicity.


\section{The exact arithmetic baseline at time zero}\label{sec:prime}

The Green identity is geometric. On the zeta slice, the safe boundary
and center can be removed exactly in favor of a convergent prime sum.
This step preserves the signed boundary information.

\begin{proposition}[Prime decomposition]\label{prop:prime}
Let
\[
 \mathcal L(X)=\int_\R K(x-X)\log|\zeta(1/2-\ii x/2)|\dd x,
\]
and
\begin{equation}\label{eq:P}
 P(X)=\sum_{n\ge2}\frac{\Lambda_{\rm vM}(n)}{\log n}
       \frac{n^{-5/2}}{1+n^{-4}}
       \cos\left(\frac X2\log n\right).
\end{equation}
Then the logarithmic integral is absolutely convergent, the prime-power
series converges absolutely and uniformly on \(\R\), and
\begin{equation}\label{eq:prime}
 J_0(X)=\mathcal L(X)-P(X)+G(X,1).
\end{equation}
For an integer \(Q\ge2\), its truncation satisfies
\begin{equation}\label{eq:prime-tail}
 |P(X)-P_Q(X)|\le\beta_Q:=\frac23Q^{-3/2}.
\end{equation}
\end{proposition}
\begin{proof}
Under \(s(z)=1/2-\ii z/2\), the heights \(0,4,8\) correspond to
real parts \(1/2,5/2,9/2\). In the xi factorization
\eqref{eq:xi}, the logarithmic moduli of \(s\), \(\pi^{-s/2}\),
\(\Gamma(s/2)\), and the constant are harmonic on a neighborhood
of the closed strip. The zero of \(s\) is at \(-\ii\), outside
the strip. Gamma has neither zeros nor poles there, and Stirling
growth makes its logarithmic modulus at most linear plus logarithmic
in horizontal height \cite{DLMFGreenPrime}. Exponential Poisson
decay therefore removes vertical-side terms and these harmonic
contributions have zero \(\calJ\).

The factor \(s-1\) has its zero at \(z=\ii\) and contributes
\emph{plus} \(G(X,1)\). This accounts for the pole of zeta; the
zeta logarithm is not treated as harmonic across \(s=1\).
Subtracting the elementary and gamma logarithms from
Lemma~\ref{lem:unit-log} also proves weighted absolute integrability
of the critical-line logarithm.

For \(\sigma>1\), the Euler product gives
\[
 \log|\zeta(\sigma-\ii x/2)|
 =\sum_{n\ge2}\frac{\Lambda_{\rm vM}(n)}{\log n}
       n^{-\sigma}\cos((x/2)\log n).
\]
The substitution \(u=\ee^{2at}\) and the beta integral yield
\[
 \widehat K(\omega)
 =\frac1{16a}\int_0^\infty
       \frac{u^{-1/2-\ii\omega/(2a)}}{1+u}\dd u
 =\frac12\operatorname{sech}(4\omega).
\]
At \(\omega=(\log n)/2\), this is \((n^2+n^{-2})^{-1}\).
The safe-boundary minus center coefficient is consequently
\[
 \frac{n^{-9/2}}{n^2+n^{-2}}-n^{-5/2}
 =-\frac{n^{-5/2}}{1+n^{-4}}.
\]
Absolute convergence justifies the interchanges and proves
\eqref{eq:prime}. Finally
\(\Lambda_{\rm vM}(p^k)/\log(p^k)=1/k\), and the coefficient
vanishes off prime powers. Thus
\[
 |P-P_Q|\le\sum_{n>Q}n^{-5/2}
          \le\int_Q^\infty t^{-5/2}\dd t=\frac23Q^{-3/2}.
\]
The integer cutoff is part of this integral-test bound.
\end{proof}

A useful pole-cancelled field is
\begin{equation}\label{eq:F}
 q(s)=\frac{s-1}{s},\qquad
 F(z)=q(s(z))\zeta(s(z))=D(z)H_0(z),\qquad
 D(z)=\frac{16\pi^{s(z)/2}}{s(z)^2\Gamma(s(z)/2)}.
\end{equation}
The multiplier \(D\) is holomorphic and zero-free on a neighborhood
of the closed strip, with harmonic logarithm of admissible growth.
Therefore \(\calJ[F]=J_0\). On the real boundary \(|q|=1\), so
\begin{equation}\label{eq:F-prime}
 J_0(X)=\int_\R K(x-X)\log|F(x)|\dd x-P(X)+G(X,1).
\end{equation}
The plus pole correction occurs once in this \emph{single-boundary}
prime representation. It is not added again to the already
pole-cancelled \emph{two-boundary} functional \(\calJ[F]\).

Positivity proves the lower comparison
\(\mathcal L(X)\ge P(X)-G(X,1)\).
The exclusion problem instead asks for an upper comparison
\[
 \mathcal L(X)<P(X)-G(X,1)+m(d,r).
\]
No such favorable bound is supplied by Proposition~\ref{prop:prime}.
No Euler product is assigned to nonzero Newman time.


\section{A square-root obstruction to uniform resolvent stability}\label{sec:conditioning}

The difficulty in \eqref{eq:ratio-error} is not removed merely by
imposing common source and zero-count bounds. The following family
makes the loss quantitative while fixing the translation scale.

Let \(H=H_0\), \(\alpha=(\log2)/2\), and
\[
 S_\alpha(u)=\frac{\Phi(u+\alpha)+\Phi(u-\alpha)}2,\qquad
 \Psi_\eta(u)=\frac{(1+\eta)\Phi(u)+(1-\eta)S_\alpha(u)}2,
\]
where \(0\le\eta\le\eta_0:=\tanh^2(\alpha/4)\).
The transforms are
\begin{equation}\label{eq:control}
 f_\eta(z)=\frac12\int_\R\Psi_\eta(u)\ee^{\ii zu}\dd u
 =H(z)M_\eta(z),\qquad
 M_\eta(z)=\cos^2(\alpha z/2)+\eta\sin^2(\alpha z/2).
\end{equation}
Every source is strictly positive, smooth and even, and
\(f_\eta(0)=H(0)\). The parameter \(\eta\) varies the source; it is
not Newman time. Each source has its own genuine heat orbit, but
\eqref{eq:control} is the time-zero factorization.

\begin{theorem}[Uniform square-root conditioning]\label{thm:conditioning}
The family \(\Psi_\eta\) has common bounds in every fixed exponentially
weighted source seminorm and common full zero-count constants. Its
relative source distance from \(\Psi_0\) is at most \(\eta\), and
each fixed safe-line logarithmic-resolvent jet changes by \(O(\eta)\),
uniformly along that line. Nevertheless,
\begin{equation}\label{eq:delta-scale}
 c\sqrt\eta\le\calJ[f_\eta](X)-\calJ[f_0](X)
                  \le C_\alpha\sqrt\eta
\end{equation}
for every \(X\in\R\), with constants independent of \(\eta,X\).
Thus no nondecreasing modulus \(o(\sqrt t)\), in particular no
uniform Lipschitz modulus, controls the Green difference from these
observations on a class containing this family.
\end{theorem}

\begin{proof}
The exact source difference is
\[
 \Psi_\eta-\Psi_0=\frac\eta2(\Phi-S_\alpha),\qquad
 |\Psi_\eta-\Psi_0|\le\eta\Psi_0.
\]
For fixed \(L,b\ge0\) and integers \(j,k\ge0\), its differentiated
weighted \(L^1\) norm is bounded by
\[
 \frac\eta2\int_\R |u|^j\ee^{Lu^2+b|u|}
            (|\Phi^{(k)}|+|S_\alpha^{(k)}|)\dd u.
\]
This integral is finite by Lemma~\ref{lem:source} and fixed translation.
The bound is independent of horizontal observation height.

Write
\[
 \ell=\frac{2\pi}{\alpha},\qquad
 x_k=\frac{(2k+1)\pi}{\alpha},\qquad
 d(\eta)=\frac2\alpha\atanh\sqrt\eta\le\frac12.
\]
At \(\eta=0\), the multiplier has double real zeros \(x_k\).
At \(\eta>0\), they are the simple zeros \(x_k\pm\ii d(\eta)\).
Any coincidence with an \(H\)-zero adds multiplicities. The fixed
lattice spacing gives a common unit-count bound and hence a common
full count after adding the canonical zeros.

The finite-factor argument of Section~\ref{sec:green}, applied
periodically to \(M_\eta\), gives
\begin{equation}\label{eq:lattice}
 \Delta_\eta(X):=\calJ[f_\eta](X)-\calJ[f_0](X)
       =\sum_{k\in\mathbb Z}G(X-x_k,d(\eta)),\qquad
 \calJ[f_0]=\calJ[H].
\end{equation}
One can justify the periodic limit directly on the cylinder: after
subtracting the Green singularities, the remaining function is
harmonic, has the correct boundary values and is bounded. Uniqueness
fixes its harmonic part. Alternatively \eqref{eq:G-decay} gives
absolute convergence of the periodized kernels. The real doubles of
\(M_0\) have zero interior contribution.

The nearest lattice point is at distance at most \(\ell/2\).
Using \eqref{eq:margin} for that point and
\eqref{eq:G-decay} for the complete lattice gives
\begin{equation}\label{eq:lattice-bounds}
 \frac{\sin a}{\cosh(a\ell/2)}\,d
 \le\Delta_\eta(X)
 \le K_G\coth(a\ell/2)\,d .
\end{equation}
The upper bound uses
\(\sup_X\sum_k\ee^{-a|X-x_k|}=\coth(a\ell/2)\).
Finally
\[
 \frac2\alpha\sqrt\eta\le d(\eta)
 \le\frac{2}{\alpha(1-\eta_0)}\sqrt\eta,
\]
which proves the two-sided Green scale.

It remains to verify the resolvent observations. With
\(R_\eta=f_\eta'/f_\eta\), \(R_*=f_0'/f_0\), and
\(B_\eta=(1+\eta)/(1-\eta)=\cosh(\alpha d)\), exact subtraction yields
\begin{equation}\label{eq:resolvent-difference}
 R_\eta-R_*=
 \frac{\alpha(B_\eta-1)\sin(\alpha z)}
 {(B_\eta+\cos(\alpha z))(1+\cos(\alpha z))}.
\end{equation}
The canonical resolvent cancels, but the multiplier's zeros remain
fully represented in \eqref{eq:lattice}. For \(y>1/2\) and
\(1\le B\le\cosh(\alpha/2)\),
\[
 |B+\cos(\alpha(x+\ii y))|\ge\cosh(\alpha y)-B,\qquad
 |\sin(\alpha(x+\ii y))|\le\cosh(\alpha y).
\]
For the first inequality, set \(t=\cos(\alpha x)\); the squared
modulus is \(t^2+2B\cosh(\alpha y)t+B^2+\sinh^2(\alpha y)\),
minimized at \(t=-1\).

Fix \(b>1\) and \(0<\varrho<b-1\), and define
\[
 C_{b,\varrho}=
 \frac{2\alpha\cosh(\alpha(b+\varrho))}
 {(1-\eta_0)
 [\cosh(\alpha(b-\varrho))-\cosh(\alpha/2)]
 [\cosh(\alpha(b-\varrho))-1]}.
\]
On every radius-\(\varrho\) disk centered at \(x+\ii b\),
\eqref{eq:resolvent-difference} is bounded by \(C_{b,\varrho}\eta\).
The actual resolvents are defined there by the zero strip.
Cauchy's estimate gives
\begin{equation}\label{eq:jet-change}
 \sup_x|\partial_z^j(R_\eta-R_*)(x+\ii b)|
 \le j!\varrho^{-j}C_{b,\varrho}\eta.
\end{equation}
Any fixed sum, maximum, or ordinary finite-dimensional norm of the
relative source distance and finitely many such jet seminorms is
therefore at most \(C_D\eta\), with weights independent of \(\eta\).
A nondecreasing \(\omega(t)=o(\sqrt t)\) would imply
\(c\sqrt\eta\le\omega(C_D\eta)=o(\sqrt\eta)\), a contradiction.
\end{proof}

The baseline \(\Psi_0\) is itself a translated-source mixture, not
theta. The theorem concerns continuity on that enlarged information
class. It is not a discontinuity statement, a universal upper
Hölder theorem, or an obstruction restricted to the literal theta
source. It leaves open \(d^2\)-scale source resolution, contours
approaching the event with their denominator costs paid, and genuinely
arithmetic signed estimates. The explicit factorial in
\eqref{eq:jet-change} also prevents interpreting fixed-order
control as a free growing-order theorem.


\section{Quadratic comparison and the short-mollifier floor}\label{sec:quadratic}

\subsection{The exact mean and the multiplier's zeros}

Fix \(0<\theta<1/2\), put \(S=T/2\), and consider
\begin{equation}\label{eq:A-class}
 A(s)=\sum_{n\le N}a_nn^{-s},\qquad
 N\le S^\theta,\quad a_1=1,\quad a_n\in\C,\quad |a_n|\le1.
\end{equation}
For each \(T\), one polynomial is used on the whole boundary and for
all centers. Set
\[
 F_A(z)=F(z)A(s(z)),\quad c_A(X)=F_A(X+4\ii),\quad
 \epsilon_A(X)^2=\int\left|\frac{F_A}{c_A(X)}-1\right|^2\dd\mu_X .
\]
The denominator is the actual center value, not a surrogate constant.

For \(\Re s>0\), the fractional-part formula gives
\begin{equation}\label{eq:fractional}
 q(s)\zeta(s)=1-(s-1)\int_1^\infty\{v\}v^{-s-1}\dd v.
\end{equation}
It follows by integrating the step-counting formula for zeta initially
on \(\Re s>1\); the displayed integral then continues the pole-cancelled
field to \(\Re s>0\). At \(s=1\) its value is one. If
\(C_A=\sum_{n\le N}|a_n|n^{-1/2}\le2\sqrt N\), then on the closed
strip \(0\le\Im z\le8\),
\(|F_A(x+\ii y)|\le C_A(8+|x|)\).
Poisson representation of its real and imaginary parts therefore gives
\begin{equation}\label{eq:mean}
 \int F_A\dd\mu_X=c_A(X).
\end{equation}
Polynomial growth is killed by the exponentially small vertical-side
harmonic measures.

At \(s_c=5/2-\ii X/2\),
\[
 2-\zeta(5/2)\le|A(s_c)|\le\zeta(5/2),\qquad
 \frac{\zeta(5)}{\zeta(5/2)}\le|\zeta(s_c)|\le\zeta(5/2),\qquad
 \frac35\le|q(s_c)|\le1.
\]
The Euler lower bound uses
\(\prod_p(1+p^{-5/2})^{-1}=\zeta(5)/\zeta(5/2)\).
Since \(\zeta(5/2)\le5/3\), the constants
\begin{equation}\label{eq:center-bounds}
 c_*=\frac35\frac{\zeta(5)}{\zeta(5/2)}(2-\zeta(5/2))>0,\qquad
 U_*=\zeta(5/2)^2
\end{equation}
satisfy \(c_*\le|c_A(X)|\le U_*\).
Expanding the square using \eqref{eq:mean} yields
\begin{equation}\label{eq:variance}
 V_A(X):=\int|F_A|^2\dd\mu_X-|c_A(X)|^2
 =|c_A(X)|^2\epsilon_A(X)^2\ge0,\qquad
 \epsilon_A(X)^2\ge V_A(X)/U_*^2.
\end{equation}

\begin{proposition}[Quadratic upper comparison]\label{prop:quadratic}
With the class \eqref{eq:A-class}, all multiplier boundary logarithms
are absolutely integrable and
\begin{equation}\label{eq:quadratic}
 0\le J_0(X)\le\calJ[F_A](X)
 =J_0(X)+\calJ[A\circ s](X)
 \le\frac12\log(1+\epsilon_A(X)^2).
\end{equation}
Zeros of \(A\), including boundary zeros, are retained.
\end{proposition}
\begin{proof}
For \(0<\delta\le1\), the bounded subharmonic function
\[
 u_\delta(z)=\frac12\log(|A(s(z))|^2+\delta^2)
\]
has upper bound \(M=\log(C_A+1)\).
Subharmonic Poisson majorization implies
\[
 0\le\int(M-u_\delta)\dd\mu_X
 \le M-u_\delta(X+4\ii)\le M-\log|A(s_c)|.
\]
Monotone convergence as \(\delta\downarrow0\) proves integrability
of the negative logarithm; its positive part is bounded. It also gives
\(\calJ[A\circ s]\ge0\). In the exact factorization
\(F_A=DH_0(A\circ s)\), the multiplier \(D\) has zero Green
functional by Section~\ref{sec:prime}. Add the logarithms and use
concavity of \(\log\):
\[
 \calJ[F_A]=\frac12\int\log|F_A/c_A|^2\dd\mu_X
 \le\frac12\log\int|F_A/c_A|^2\dd\mu_X
 =\frac12\log(1+\epsilon_A^2).
\]
No zero of \(A\) was divided away, and the zeta pole was already
cancelled in \(F_A\).
\end{proof}

\subsection{The imported mean-value input and an elementary first moment}

We use the unconditional lower-bound part of
Radziwi{\l}{\l}'s Proposition B, attributed there to Soundararajan
\cite{Radziwill2012MollifyingV3}. For fixed \(\theta<1/2\),
normalization \(a_1=1\), and the stated common coefficient-growth
bounds, it gives
\begin{equation}\label{eq:primary-moment}
 \frac1S\int_S^{2S}
 |1-\zeta(1/2+\ii t)A(1/2+\ii t)|^2\dd t
 \ge\frac1\theta-o_\theta(1).
\end{equation}
This is an imported mean-value theorem, not proved here. We use the
pinned version 3 and its complex-coefficient statement, not its
RH-conditional results.

For \eqref{eq:A-class}, the lower remainder is uniform in \(A\).
Pad shorter polynomials with zeros to \(\lfloor S^\theta\rfloor\).
All coefficient-growth constants are at most one. If a uniform lower
remainder failed, offending arrays along a sequence of heights,
extended by \(A=1\) at the other heights, would form an admissible
family contradicting \eqref{eq:primary-moment}. This does not assert
an additive uniform asymptotic for arbitrarily large quadratic forms.

To convert the error moment to a square mean, an additional first
moment is essential. Put \(B(s)=\zeta(s)A(s)\).
Shift its rectangle from \(\Re s=1/2\) to \(\Re s=2\) between
imaginary heights \(S\) and \(2S\). The zeta pole is outside.
For the integer \(M=\lceil2S\rceil\), the exact truncation
\[
 \zeta(s)=\sum_{n\le M}n^{-s}+\frac{M^{1-s}}{s-1}
         -s\int_M^\infty\{v\}v^{-s-1}\dd v
\]
gives \(|\zeta(s)|=O(\sqrt S)\) on both horizontal edges,
uniformly for \(1/2\le\Re s\le2\). Since \(|A(s)|\le2\sqrt N\),
the two edges cost \(O(\sqrt{SN})\).
On \(\Re s=2\), absolute convergence gives
\[
 B(s)=\sum_{n\ge1}b_nn^{-s},\qquad
 b_n=\sum_{\substack{d\mid n\\d\le N}}a_d,\quad b_1=1,\quad
 |b_n|\le\tau(n).
\]
The nonconstant mean has absolute size at most
\(2S^{-1}\sum_{n\ge2}\tau(n)n^{-2}/\log n=O(S^{-1})\).
Consequently
\begin{equation}\label{eq:first-moment}
 \frac1S\int_S^{2S}B(1/2+\ii t)\dd t
 =1+O_\theta(S^{(\theta-1)/2}).
\end{equation}
Expanding \(|1-B|^2\) in \eqref{eq:primary-moment} now gives
\begin{equation}\label{eq:square-mean}
 \frac1S\int_S^{2S}|B(1/2+\ii t)|^2\dd t
 \ge1+\frac1\theta-o_\theta(1),
\end{equation}
uniformly in the declared coefficient class.

\subsection{Safe-line subtraction and the padded convolution}

For any interval \([u,u+S]\), absolute convergence gives
\begin{equation}\label{eq:safe-mean}
 \frac1S\int_u^{u+S}|B(5/2+\ii t)|^2\dd t
 =\sum_{n\ge1}|b_n|^2n^{-5}+O(S^{-1}),
\end{equation}
uniformly in \(u,A\).
Indeed the absolute off-diagonal error is at most
\[
 \frac4S\sum_{m<n}
 \frac{\tau(m)\tau(n)}{(mn)^{5/2}\log(n/m)}
 \le\frac{16}S\sum_{m\ge1}\frac{H_m}{m^3}<\infty,
 \quad H_m=\sum_{j=1}^m\frac1j.
\]
Here \(\log(n/m)\ge(n-m)/n\), \(\tau(n)\le2\sqrt n\), and
\(\sum_{k\ge1}1/[k(m+k)]=H_m/m\).
The diagonal is bounded by
\begin{equation}\label{eq:diagonal}
 \sum_{n\ge1}\tau(n)^2n^{-5}
 \le1+4\sum_{n\ge2}n^{-4}
 \le1+4\left(\frac1{16}+\int_2^\infty v^{-4}\dd v\right)
 =\frac{17}{12}.
\end{equation}

Let
\begin{equation}\label{eq:padded}
 L_T=a^{-1}\log T,\qquad I_T=[T-L_T,2T+L_T].
\end{equation}
Since \(|q(5/2+\ii t)|\le1\), substitution \(X=2t\), with
coefficient conjugation for the negative imaginary direction, gives
\[
 \frac1T\int_T^{2T}|c_A(X)|^2\dd X\le17/12+O(T^{-1}).
\]
The two extra intervals have total length \(2L_T\), and
\(|c_A|^2\le U_*^2\). Normalizing by the actual length therefore yields
\begin{equation}\label{eq:padded-center}
 \frac1{|I_T|}\int_{I_T}|c_A(X)|^2\dd X\le17/12+o(1).
\end{equation}

\begin{theorem}[Uniform quadratic-residual floor]\label{thm:floor}
For every fixed \(0<\theta<1/2\), the class \eqref{eq:A-class}
and interval \eqref{eq:padded} satisfy
\begin{equation}\label{eq:floor}
 \liminf_{T\to\infty}\inf_A
 \frac1{|I_T|}\int_{I_T}\epsilon_A(X)^2\dd X
 \ge\frac{1/(2\theta)-11/12}{\zeta(5/2)^4}
 >\frac1{12\zeta(5/2)^4}.
\end{equation}
In particular, for all sufficiently large \(T\), uniformly in \(A\),
the average exceeds \(1/[24\zeta(5/2)^4]\).
No numerical onset is asserted.
\end{theorem}
\begin{proof}
Integrate the complete nonnegative boundary square in
\eqref{eq:variance} over \(I_T\), using Tonelli.
Retain only the lower boundary and \(x\in[T,2T]\).
For every such \(x\), the interval
\([x-L_T,x+L_T]\) lies in \(I_T\), and hence
\[
 \int_{I_T}K(x-X)\dd X\ge k_{L_T},\qquad
 k_L=\int_{-L}^LK(u)\dd u=\frac1\pi\arctan(\sinh(aL))
 \ge\frac12-\frac{\ee^{-aL}}{4a}.
\]
Thus \(k_{L_T}\ge1/2-O(T^{-1})\).
On the lower boundary \(|q|=1\).
The conjugate-coefficient polynomial
\(\widetilde A(s)=\sum\overline{a_n}n^{-s}\) is still admissible.
Applying \eqref{eq:square-mean} to it and substituting \(x=2t\)
gives
\[
 \int_T^{2T}|F_A(x)|^2\dd x
 =2\int_S^{2S}|\zeta(1/2+\ii t)\widetilde A(1/2+\ii t)|^2\dd t
 \ge T(1+1/\theta-o_\theta(1)).
\]
Subtract \eqref{eq:padded-center}. Since \(T/|I_T|\to1\),
\[
 \frac1{|I_T|}\int_{I_T}V_A(X)\dd X
 \ge\frac12(1+1/\theta)-\frac{17}{12}-o_\theta(1)
 =\frac1{2\theta}-\frac{11}{12}-o_\theta(1).
\]
All lower errors are common to the coefficient class. Divide by
\(U_*^2=\zeta(5/2)^4\) using \eqref{eq:variance}, then take the
infimum and liminf. The strict assumption \(\theta<1/2\) gives
the final claims.
\end{proof}

\subsection{The real mixed term and its corrected verification}

For clarity we include the elementary verification used in the branch's
assessment \cite{OGCorrection,OGAssessment}. It corrects an intermediate
source attribution and does not replace the imported mean-value
theorem. With \(N=\lfloor S^\theta\rfloor\), define
\[
 w_\ell=\frac{\varphi(\ell)}{\ell^2},\quad
 y_\ell=\sum_{d\le N/\ell}\frac{a_{d\ell}}d,\quad
 G_N=\sum_{\ell\le N}\frac{\mu(\ell)^2}{\varphi(\ell)},\quad
 z_\ell=\frac{\mu(\ell)\ell}{G_N\varphi(\ell)},
\]
\[
 e_\ell=y_\ell-z_\ell,\qquad
 E_N=\sum_{\ell\le N}w_\ell|e_\ell|^2.
\]
Möbius inversion and \(a_1=1\) imply
\begin{equation}\label{eq:orthogonality}
 \sum_{\ell\le N}\frac{\mu(\ell)e_\ell}{\ell}=0.
\end{equation}
The relevant \emph{real} mixed form is
\[
 \mathcal C_N=
 2\Re\sum_{p^k\ell\le N}\frac{\log p}{p^k}w_\ell
       \{z_\ell\overline{e_{p^k\ell}}+z_{p^k\ell}e_\ell\}.
\]
Call the two directed sums \(U_N,V_N\).
Reindexing the first and using
\(\sum_{p^k\mid m}\mu(m/p^k)\log p=-\mu(m)\log m\) gives
\[
 U_N=-\frac1{G_N}\sum_{m\le N}
             \frac{\mu(m)\log m}{m}\overline{e_m}.
\]
In the second, \(\mu(p^k\ell)\ne0\) requires \(k=1\), \(p\nmid\ell\).
Thus the exact expression is
\[
 V_N=-\frac1{G_N}\sum_{\ell\le N}\frac{\mu(\ell)e_\ell}{\ell}
         P_\ell(N/\ell),\qquad
 P_\ell(x)=\sum_{\substack{p\le x\\p\nmid\ell}}\frac{\log p}{p-1}.
\]
The classical prime-sum estimate used in the primary argument gives
\[
 P_\ell(N/\ell)=\log(N/\ell)+r_\ell,\qquad
 |r_\ell|\ll\log\log(3N).
\]
Its uniformity follows by splitting excluded prime divisors at
\(Y=\log(3N)\): the small ones cost \(O(\log Y)\), and the large
ones cost at most \(2\log N/Y\).
The unrestricted sum is \(\log x+O(1)\); replacing
\(1/(p-1)\) by \(1/p\) costs a convergent series.

After \eqref{eq:orthogonality}, the logarithmic main terms of \(U_N\)
and \(V_N\) cancel \emph{only after taking real parts}. Therefore
\[
 \mathcal C_N=-\frac2{G_N}\Re
 \sum_{\ell\le N}\frac{\mu(\ell)r_\ell e_\ell}{\ell}.
\]
Cauchy--Schwarz and the classical \(G_N\sim\log N\) give
\begin{equation}\label{eq:real-mixed}
 |\mathcal C_N|
 \le\frac2{G_N}
 \left(\sum_{\ell\le N}\frac{\mu(\ell)^2r_\ell^2}{\varphi(\ell)}\right)^{1/2}
 \sqrt{E_N}
 \ll\frac{\log\log(3N)}{\sqrt{\log N}}\sqrt{E_N}.
\end{equation}
For a fixed \(\nu>0\), if the last coefficient is \(b_N\), then
\[
 b_N\sqrt{E_N}
 \le\nu\log S\,E_N+\frac{b_N^2}{4\nu\log S}.
\]
The scalar remainder tends to zero uniformly in the coefficients.
The two elementary prime-sum asymptotics are classical inputs already
used in the source argument. Equation~\eqref{eq:real-mixed} is the
direct real calculation recorded by this project; it is not quoted
as the literal printed formula of Lemma 13 in version 3, which displays
\(E_N\) rather than its square root. Nor does it bound the modulus
of the unprojected complex mixed sum.

\subsection{What the floor excludes}

By Corollary~\ref{cor:margin} and Proposition~\ref{prop:quadratic},
a fully acquired inequality
\(\epsilon_A(X)^2<\exp(2m(d,r))-1\) would exclude the specified
nonreal zero. The positive average floor prevents these sufficient
residual certificates from tending uniformly to zero throughout
\(I_T\) using one polynomial in \eqref{eq:A-class}.
It gives no lower bound for \(J_0\), because the upper comparison
\eqref{eq:quadratic} cannot be reversed. It supplies neither failure
at every center nor a positive density of failing centers.
Longer polynomials, \(\theta\ge1/2\), adaptive center choices,
unpadded averages and signed logarithmic estimates are outside
the theorem.


\section{Complete source acquisition for the quadratic comparison}\label{sec:acquisition}

The preceding floor concerns the exact field. To state honestly what
a finite implementation of its sufficient comparison would require,
we retain the complete acquisition estimate from
\cite{OGMollifier,OGCorrection}. This is an analytical error bound;
no numerical certificate has been run.

\subsection{Arithmetic and radial source errors}

Let \(\Phi_{N_{\rm src}}\) be the first \(N_{\rm src}\) summands of
\eqref{eq:source} on \(u\ge0\), reflected evenly, and put
\[
 \widehat H(z)=\frac12\int_{-V}^{V}
          \Phi_{N_{\rm src}}(u)\ee^{\ii zu}\dd u.
\]
With \(B_0=5C_2(0)\), define
\[
 \mathcal Q=\int_\R\ee^{9|u|}\rho(u)^2\ee^{-\rho(u)}\dd u,\qquad
 \mathcal Q^{>V}=\int_{|u|>V}\ee^{9|u|}
                      \rho(u)^2\ee^{-\rho(u)}\dd u.
\]
Both integrals are finite by Section~\ref{sec:source}. Splitting the
difference into the arithmetic tail and the exact radial tail gives
\begin{equation}\label{eq:delta-H}
 \sup_{|\Im z|\le8}|H_0(z)-\widehat H(z)|
 \le\delta_H:=\frac12\{5C_2(N_{\rm src})\mathcal Q+
                          B_0\mathcal Q^{>V}\}.
\end{equation}
The terms tend to zero separately. Only the undifferentiated source
is used, so no smooth matching of the finite prefix at zero is assumed.

For \(|x-X|\le R\), \(0\le\Im z\le8\), the reciprocal-gamma
product \cite{DLMFReciprocalGammaOC} gives the compact multiplier bound
\begin{equation}\label{eq:multiplier}
 |D(z)A(s(z))|\le\mathcal B_A(X,R)
 :=16\pi^{9/4}C_A
 \exp\left\{\frac94+\frac{\pi^2[81+(|X|+R)^2]}{192}\right\}.
\end{equation}
To verify it, write \(w=s/2\) in
\[
 \Gamma(w)^{-1}=w\ee^{\gamma w}
        \prod_{n\ge1}(1+w/n)\ee^{-w/n}.
\]
The inequality
\(\log|1+w/n|\le\Re w/n+|w|^2/(2n^2)\) yields
\[
 |\Gamma(w)^{-1}|
 \le|w|\exp\{\gamma\Re w+\pi^2|w|^2/12\}.
\]
Now \(1/2\le\Re s\le9/2\), \(|s|\ge1/2\),
\(|s|^2\le[81+(|X|+R)^2]/4\), and \(\gamma<1\).
Insert these bounds into \eqref{eq:F} and use \(|A|\le C_A\).
With \(\widehat F_A=DA\widehat H\), the compact field error is at
most \(\mathcal B_A\delta_H\).
This bound is not used on the remote infinite boundary.

\subsection{Both exact boundary tails}

Let \(\mu_{X,R}\) be the restriction of \(\mu_X\) to \(|x-X|\le R\).
Equation~\eqref{eq:fractional} gives
\[
 |F_A(x)|\le C_A(2+|x|),\qquad
 |F_A(x+8\ii)|\le\zeta(9/2)^2.
\]
Writing \(B_X=2+|X|\), one obtains
\begin{align}\label{eq:quadratic-tail}
 \norm{F_A/c_A-1}_{L^2(\mu_X-\mu_{X,R})}^2
 &\le\tau_A(X,R),\nonumber\\
 \tau_A(X,R)
 &=\frac{\ee^{-aR}}{2a}\left\{
 \frac{C_A^2}{c_*^2}
 \left[(B_X+R)^2+\frac{2(B_X+R)}a+\frac2{a^2}\right]
 +\frac{\zeta(9/2)^4}{c_*^2}+2\right\}.
\end{align}
Indeed \(|v-1|^2\le2|v|^2+2\), and
\(K(u)\le\ee^{-a|u|}/8\).
Integrate the lower-boundary quadratic polynomial on both spatial
tails and then the upper-boundary constant on both tails.
The upper square requires the fourth power of \(\zeta(9/2)\).
All tails here belong to the exact field.

\subsection{Center acquisition and denominator propagation}

For an integer \(M_c\ge2\), the safe Euler logarithm has error
\[
 \left|\log\zeta(s_c)-\sum_{p^k\le M_c}\frac{p^{-ks_c}}k\right|
 \le\beta_{M_c}=\frac23M_c^{-3/2}.
\]
Keep the finite polynomial \(A(s_c)\) exactly and set
\[
 \widetilde c=q(s_c)A(s_c)
       \exp\left(\sum_{p^k\le M_c}\frac{p^{-ks_c}}k\right),\qquad
 |c_A-\widetilde c|\le
 \eta:=U_*(\ee^{\beta_{M_c}}-1).
\]
The modulus of the finite-log exponential is at most \(\zeta(5/2)\).
Future arithmetic evaluation error is added to \(\eta\).
No logarithm of \(A\) is used.

Let
\[
 \ell_c=|\widetilde c|-\eta>0,\qquad
 \mu_R=\mu_{X,R}(\partial{\rm strip})
       =\frac2\pi\arctan(\sinh(aR)),\qquad
 A_R=\norm{\widehat F_A/\widetilde c-1}_{L^2(\mu_{X,R})}.
\]
Then
\begin{equation}\label{eq:full-quadratic-acquisition}
 \epsilon_A(X)\le A_R+
 \frac{\eta}{\ell_c}(A_R+\sqrt{\mu_R})
 +\frac{\mathcal B_A\delta_H}{c_*}+\sqrt{\tau_A(X,R)}.
\end{equation}
For the proof, first split the complete \(L^2\) norm into compact
and remote pieces using the triangle inequality. On the compact
piece, add and subtract \(\widehat F_A/c_A\) and
\(\widehat F_A/\widetilde c\). The field error is at most
\(\mathcal B_A\delta_H/c_*\), since the restricted mass is at most one.
The denominator error is
\[
 \frac{|c_A-\widetilde c|}{|c_A|}
     \norm{\widehat F_A/\widetilde c}
 \le\frac{\eta}{\ell_c}(A_R+\sqrt{\mu_R}),
\]
which proves \eqref{eq:full-quadratic-acquisition}.

For each fixed \(X,A\) and a positive acquisition budget \(h\),
first choose \(R\) to make \(\sqrt{\tau_A}<h/3\).
For this \(R\), choose \(N_{\rm src},V\) so that
\(\mathcal B_A\delta_H/c_*<h/3\).
Finally choose the integer center cutoff and arithmetic precision so that
\(\eta<c_*/4\) and
\(2\eta\norm{\widehat F_A}_{L^2(\mu_{X,R})}/c_*^2<h/3\).
The exact denominator difference is then bounded by
\[
 \frac{\eta\norm{\widehat F_A}}{|c_A||\widetilde c|}
 \le\frac{4\eta\norm{\widehat F_A}}{3c_*^2}<h/3.
\]
This sharper expression, rather than requiring the conservative
surrogate in \eqref{eq:full-quadratic-acquisition} to meet the same
allocation, proves the ordered finite-acquisition assertion.
If the norm vanishes, the denominator error vanishes.
Every tuple can be treated this way, including a prescribed sequence
of growing heights or shrinking defect distances. This supplies no
common cutoff, affordable complexity bound, small residual value,
or evaluated exclusion.


\section{One-sided boundary-logarithm approximation through zeros}\label{sec:boundary-log}

The quadratic floor concerns a sufficient positive surrogate. We now
retain the measured logarithm itself. The local theorem below controls
how much an absolute field error can raise its integral, even at
multiple real zeros. Its starting data are a nonzero anchor away
from the real interval and an upper bound on a larger disk.

\subsection{A one-dimensional sublevel theorem}

Fix the nested radii and constants
\begin{equation}\label{eq:disk-constants}
 R_0=\frac92,\quad \rho_0=\frac{13}3,\quad r_0=\frac{25}6,\quad
 B=R_0+r_0=\frac{26}3,\quad
 \kappa=\log\frac{27}{26},\quad
 \frac{\rho_0+r_0}{\rho_0-r_0}=51 .
\end{equation}

\begin{theorem}[Line sublevels and a logarithmic upper envelope]\label{thm:local-log}
Let \(f\) be holomorphic on a neighborhood of the closed disk
\(\overline{D(0,R_0)}\), with \(|f|\le1\) there and
\(|f(0)|\ge m>0\). Put
\begin{equation}\label{eq:n-C}
 A=\log(1/m),\qquad
 n=\max\{2,\lceil A/\kappa\rceil\},\qquad
 C_f=2Bn\ee^{51A/n}.
\end{equation}
Let \(\Gamma\) be a horizontal segment of length at most one inside
\(\overline{D(0,r_0)}\). Then, for every \(t>0\),
\begin{equation}\label{eq:sublevel}
 |\{z\in\Gamma:|f(z)|<t\}|\le C_ft^{1/n}.
\end{equation}
If \(g\) is measurable on \(\Gamma\) and
\(|f-g|\le\varepsilon\), with \(\varepsilon>0\), then
\begin{equation}\label{eq:local-envelope}
 0\le\int_\Gamma
       [\log(|g|+\varepsilon)-\log|f|]\dd\ell
 \le\frac{C_fn^2}{n-1}(2\varepsilon)^{1/n}.
\end{equation}
No lower bound for \(|f|\) on \(\Gamma\), simplicity or separation
of its zeros, or nonvanishing or holomorphy of \(g\), is assumed.
\end{theorem}

\begin{proof}
Jensen's formula \cite{Tao2020LocalZeros} bounds the number \(k\)
of zeros inside the open \(\rho_0\)-disk, including multiplicity,
by \(k\kappa\le A\). Remove exactly these zeros with disk Blaschke
factors:
\[
 b_j(z)=\frac{R_0(z-a_j)}{R_0^2-\overline{a_j}z},\qquad
 h(z)=f(z)\Big/\prod_{j=1}^k b_j(z).
\]
Fill the removable singularities. Since every factor has boundary
modulus one, the maximum principle gives \(|h|\le1\) on the outer
disk. It is zero-free on the open \(\rho_0\)-disk and
\(|h(0)|\ge m\). Harnack's inequality for the nonnegative harmonic
function \(-\log|h|\) gives
\(|h(z)|\ge\ee^{-51A}\) when \(|z|\le r_0\)
\cite{Tao2016Harnack}. If \(h\) has a zero on the \(\rho_0\)-circle,
apply Harnack on smaller disks and pass to the limit.
The elementary denominator bound for each Blaschke factor gives
\begin{equation}\label{eq:factor-lower}
 |f(z)|\ge\ee^{-51A}
       \prod_{j=1}^k\frac{|z-a_j|}{B}\qquad(|z|\le r_0).
\end{equation}
Every distance ratio is less than one.

Let \(q=\ee^{51A}t\). If \(q<1\) and \(k=0\), the sublevel set
is empty. Otherwise, when \(|f(z)|<t\), at least one distance
is below \(Bq^{1/k}\). The set is covered by \(k\) disks of this
radius. A disk of radius \(b\) meets a line in length at most \(2b\),
so the total length is at most
\[
 2Bkq^{1/k}\le2Bnq^{1/n}=C_ft^{1/n}.
\]
For \(q\ge1\), this bound exceeds the segment length and is
automatic. This proves \eqref{eq:sublevel} in one dimension.
An area estimate would not suffice for this trace argument.

The triangle inequality gives, off the finitely many segment zeros,
\[
 0\le\log(|g|+\varepsilon)-\log|f|
 \le\log(1+2\varepsilon/|f|).
\]
Write \(h_*=2\varepsilon\). Tonelli and the layer-cake identity give
\[
 \int_\Gamma\log(1+h_*/|f|)\dd\ell
 =\int_0^\infty
 \frac{h_*}{t(t+h_*)}|\{|f|<t\}\cap\Gamma|\dd t.
\]
On \(0<t<h_*\), bound the kernel by \(1/t\); on \(t>h_*\),
bound it by \(h_*/t^2\). Equation~\eqref{eq:sublevel} then gives
the two bounds \(C_fn h_*^{1/n}\) and
\(C_fn h_*^{1/n}/(n-1)\). Their sum is
\eqref{eq:local-envelope}. The same sublevel integral proves local
integrability of \(\log^-|f|\). The choice \(n\ge2\) makes the
second integral finite and avoids the endpoint logarithm at \(n=1\).
\end{proof}

An immediate consequence, for \(M\ge1\), is
\begin{equation}\label{eq:absolute-log-local}
 \int_\Gamma|\log|Mf||\dd\ell\le\log M+C_fn.
\end{equation}
Indeed \(|\log|Mf||\le\log M+\log^-|f|\), and integrating
\eqref{eq:sublevel} over \(0<t<1\) with measure \(\dd t/t\)
bounds the latter integral by \(C_fn\).

\subsection{The canonical anchor and explicit height dependence}

For \(v\in\R\), center these disks at \(c_v=v+4\ii\), and take
the segment \(I_v=[v-1/2,v+1/2]\). Its maximum distance from
the center is \(\sqrt{65}/2<r_0\).
Throughout the outer disk, \(1/4\le\Re s\le19/4\).
The pole-cancelled fractional-part identity
\eqref{eq:fractional} gives
\[
 |F(z)|\le2|v|+25.
\]
In detail \(|s-1|\le |v|/2+6\), while the integral in that
identity is bounded by \(1/\Re s\le4\). At the center,
the absolutely convergent Euler product gives
\begin{equation}\label{eq:canonical-anchor}
 |F(c_v)|\ge c_0:=\frac35\frac{\zeta(5)}{\zeta(5/2)}>0.
\end{equation}
This safe-line estimate is unconditional and supplies no lower
modulus on the real segment.

Define
\begin{equation}\label{eq:local-canonical-constants}
 M_v=2|v|+26,\quad A_v=\log(M_v/c_0),\quad
 n_v=\max\{2,\lceil A_v/\kappa\rceil\},\quad
 C_v=2Bn_v\ee^{51A_v/n_v}.
\end{equation}
Apply Theorem~\ref{thm:local-log} to \(F/M_v\).

\begin{corollary}[Canonical local acquisition]\label{cor:canonical-log}
If \(\sup_{I_v}|F-\widehat F|\le\delta\), with \(\delta>0\), then
\begin{equation}\label{eq:canonical-envelope}
 0\le\int_{I_v}
 [\log(|\widehat F|+\delta)-\log|F|]\dd x
 \le\mathcal E_v(\delta)
 :=\frac{C_vn_v^2}{n_v-1}
       \left(\frac{2\delta}{M_v}\right)^{1/n_v}.
\end{equation}
Also
\begin{equation}\label{eq:canonical-absolute-log}
 \int_{I_v}|\log|F(x)||\dd x\le\log M_v+C_vn_v.
\end{equation}
\end{corollary}
\begin{proof}
Use \(f=F/M_v\), \(g=\widehat F/M_v\),
\(\varepsilon=\delta/M_v\). Logarithmic constants cancel
exactly in the difference. Equation~\eqref{eq:absolute-log-local}
gives the second assertion.
\end{proof}

Because \(A_v/n_v\le\kappa\), the prefactor in
\eqref{eq:canonical-envelope} is \(O(\log^2(2+|v|))\),
whereas \(n_v\) is of order \(\log(2+|v|)\). To make the local
error at most \(h>0\), it suffices to take
\begin{equation}\label{eq:precision}
 \delta\le\frac{M_v}{2}
 \left(\frac{h(n_v-1)}{C_vn_v^2}\right)^{n_v}.
\end{equation}
This is finite precision for each fixed interval, not a uniform
positive Hölder exponent or a useful complexity theorem.

The local double-zero model illustrates the relevant scale without
any numerical approximation:
\begin{equation}\label{eq:double-log}
 \int_\R[\log(x^2+\delta)-\log x^2]\dd x=2\pi\sqrt\delta.
\end{equation}
Substitute \(x=\sqrt\delta\,u\) and integrate
\(\log(1+u^{-2})\) by parts on the positive half-line; its integral
is \(2\int_0^\infty(1+u^2)^{-1}\dd u=\pi\).
Thus a real zero produces no fixed error floor, although a Lipschitz
error bound fails even in this simple model. This is an algebraic
scale calculation, not a canonical collision control, and the
conservative exponent in \eqref{eq:canonical-envelope} is not
claimed sharp.

The approximant in this theorem may have zeros, and its logarithmic
integral is retained. This is the precise distinction from a
zero-free harmonic-baseline replacement. Likewise no small
quadratic-residual hypothesis from Section~\ref{sec:quadratic}
is used. The theorem measures signed information already present
in the exact field; it does not establish its favorable value.


\section{Scope, evidence and outstanding tasks}\label{sec:tasks}

\subsection{Control families and the meaning of the results}

The complete Green identity is an exact representation; the stability
and residual floors are quantitative scoped obstructions; and
Corollary~\ref{cor:canonical-log} is an explicit local acquisition
estimate. These roles must not be interchanged. In particular a
lower bound for a sufficient quadratic residual is not a lower
bound for \(J_0\).

The collision controls summarized in \cite{OGControls} give the
following analytical screen. No new control computation is required
for these premise comparisons.

\begin{center}\small
\begin{tabularx}{\textwidth}{@{}>{\raggedright\arraybackslash}p{0.24\textwidth}Y@{}}
\toprule
Control family & Premise not supplied for the canonical application\\
\midrule
Euler-factor and translated-source families &
The extra multiplier changes the literal theta/xi field. Canonical
prime subtraction and Euler-anchor constants cannot simply be assigned
to the altered function. The fixed-prime baseline in
Section~\ref{sec:conditioning} is itself not theta.\\
Weighted-square, finite-prefix and finite-prime-orbit controls &
Matched finite data do not supply the complete normalized theta
coefficients together with full reflection.\\
Reflected super-Gaussian control &
Reflection and positive Gram forms do not supply the integer-square
theta expansion or the unchanged xi/zeta identity.\\
Smooth shape and positive-primitive controls &
Their specified shape, endpoint and asymptotic data do not supply the
complete canonical arithmetic source.\\
Atomic, finite-profile and variable-readout examples &
Only their actual stated source and observation premises may be used;
the full canonical source/profile/cosine-readout conjunction is absent.\\
\bottomrule
\end{tabularx}
\end{center}

The generic local theorem can apply to an admissible control with its
own anchor and growth constants. If that control has an off-axis zero
in the relevant domain, accurate acquisition retains its positive
Green signal and cannot pass the corresponding strict exclusion
threshold. Failure of the literal theta identity does not imply
failure of every inequality derived from it. A new canonical sign
estimate would have to identify its quantitative distinguishing
premise.

\subsection{Evidence and assistance}

This draft was prepared for Edward Baker with substantial OpenAI Codex
LLM assistance in derivation, exposition, source checking and production.
The current session's recorded history supplies no model or
reasoning-effort configuration; none is inferred. The source/resolvent
conditioning and short-mollifier transfer retain their internal
assessment with the additive correction
\cite{OGCorrection,OGAssessment}. The local boundary-log theorem
retains proof-author checking \cite{OGBoundaryLog}.
The parent Green representation and prime reduction retain their
existing scoped internal assessments
\cite{LocalizedMetricsAssessment,PrimeBaselineAssessment}.
The external moment theorem is an explicit imported input.
There has been no independent specialist human review or formal
verification. Manuscript integration does not strengthen these
evidence levels.

No numerical experiment is needed for any proof in this draft.
Future computations can support feasibility, error-budget evaluation
or a declared finite certificate, but cannot supply evidence for an
all-height RH-equivalent statement merely by checking finite cases.

\subsection{Outstanding analytical tasks}

\paragraph{1. Complete the direct logarithmic assembly.}
The existing frozen execution task \cite{OGAssembly} asks for a
finite-source upper approximation to \(J_0(X)\).
For integer \(R\ge1\), \(Q\ge2\), and
\(\sup_{|x-X|\le R}|F-\widehat F|\le\delta\), its proposed form is
\begin{align}\label{eq:proposed-U}
 U_{R,Q,\delta}(X)
 ={}&\int_{X-R}^{X+R}K(x-X)\log(|\widehat F(x)|+\delta)\dd x
 \nonumber\\
 &+\mathcal T(X,R)-P_Q(X)+\beta_Q+G(X,1).
\end{align}
Here \(\widehat F=D\widehat H\) uses the unchanged source,
and \(\mathcal T\) must bound the absolute omitted logarithmic tail
of the exact \(F\). The finite compact error sum proposed from
\eqref{eq:canonical-envelope} is
\[
 \mathcal E_{X,R}(\delta)=
 \sum_{k=0}^{R-1}
 \frac{\mathcal E_{X+k+1/2}(\delta)+
       \mathcal E_{X-k-1/2}(\delta)}{16\cosh(ak)}.
\]
The outstanding complete statement is
\begin{equation}\label{eq:proposed-assembly}
 0\le U_{R,Q,\delta}(X)-J_0(X)
 \le\mathcal E_{X,R}(\delta)+2\mathcal T(X,R)+2\beta_Q.
\end{equation}
Equations~\eqref{eq:proposed-U}--\eqref{eq:proposed-assembly}
are a specified \emph{target}, not a theorem proved by this manuscript.
The assembly must derive the full explicit log tail, verify the
compact multiplier/source payment, and justify the ordered finite
choices of \(R,Q,\delta,N_{\rm src},V\).
Both factors of two are intentional: replacing a signed tail by an
upper allowance can cost twice its absolute bound, and the same is
true for prime truncation. The actual signed compact integral must
be retained.

\paragraph{2. Assess the assembled claim and its local input.}
A claim-level assessment should recheck the line-sublevel theorem,
its canonical application, all source and logarithmic tails,
the single plus pole correction, and evaluation-error propagation.
The local result currently has proof-author checking; the complete
assembly has not yet been executed. Changing the model alone would
not supply specialist independence.

\paragraph{3. Acquire a favorable canonical value.}
Even after \eqref{eq:proposed-assembly}, one needs a directed evaluated
upper value satisfying
\[
 U\le\overline U\le U+\eta_{\rm eval},\qquad
 \overline U<m(d,r).
\]
Alternatively, a source-specific analytical estimate could establish
the same strict upper comparison directly.
An approximation error tending to zero does not force this value
inequality. The missing information is a quantitative bound on the
signed critical-line logarithmic average relative to its exact prime
baseline, including the exceptional centers relevant to a proposed
zero.

\paragraph{4. Resolve uniformity and practical precision.}
The present local exponent decreases with height; the compact
source multiplier in \eqref{eq:multiplier} is deliberately coarse.
One must separate finite acquisition for each prescribed tuple
from common estimates as \(|X|\to\infty\) or \(d\downarrow0\).
Sharpening constants may improve a practical diagnostic but does not
establish canonical smallness. Any new resolvent or contour route
must pay its denominators, residues, harmonic data and tails.

\paragraph{5. Use numerical evidence only for a declared finite question.}
After the analytical assembly and a finite error budget are fixed,
a bounded diagnostic may measure cost and certified accuracy against
a comparable direct method. It should have specified cases and a
stopping rule. No pilot, all-height extrapolation, favorable canonical
margin or new numerical claim is included in the present draft.


\begingroup\footnotesize\raggedright
\bibliographystyle{plainurl}
\bibliography{green_resolvent_manuscript_sources}
\endgroup
\end{document}

